\begin{lemma}
\label{lemma-base-change-scheme-theoretic-image}
Let $S$ be a quasi-compact and quasi-separated scheme.
Let $f : X \to Y$ be a morphism of schemes over $S$ with
both $X$ and $Y$ of finite presentation over $S$.
Then there exists a $t \geq 0$ and closed subschemes
$$
S \supset S_0 \supset S_1 \supset \ldots \supset S_t = \emptyset
$$
with the following properties:
\begin{enumerate}
\item $S_i \to S$ is defined by a finite type ideal sheaf,
\item $S_0 \subset S$ is a thickening, and
\item with $T_i = S_i \setminus S_{i + 1}$ and $f_i$ the base
change of $f$ to $T_i$ we have:
formation of the scheme theoretic image of $f_i/T_i$
commutes with arbitrary base change (see discussion above the lemma).
\end{enumerate}
\end{lemma}

\begin{proof}
We can find a commutative diagram
$$
\xymatrix{
X \ar[d] \ar[r] & Y \ar[d] \ar[r] & S \ar[d] \\
U \ar[r] & V \ar[r] & W
}
$$
with cartesian squares such that $U$, $V$, $W$
are of finite type over $\mathbf{Z}$. Namely, first write $S$
as a cofiltered limit of finite type schemes over $\mathbf{Z}$
with affine transition morphisms
using Limits, Proposition \ref{limits-proposition-approximate}
and then descend the morphism $X \to Y$ using
Limits, Lemma \ref{limits-lemma-descend-finite-presentation}.
This reduces us to the case discussed in the next paragraph.

\medskip\noindent
Assume $S$ is Noetherian.  In this case every quasi-coherent ideal
is of finite type, hence we do not have to check the condition that
$S_i$ is cut out by a finite type ideal. Set $S_0 = S_{red}$ equal
to the reduction of $S$. Let $\eta \in S_0$ be a generic point
of an irreducible component of $S_0$. By Noetherian induction on
the underlying topological space of $S_0$, we may assume the result
holds for any closed subscheme of $S_0$ not containing $\eta$.
Thus it suffices to show that there exists an open neighbourhood
$U_0 \subset S_0$ such that the base change $f_0$ of $f$ to $U_0$
has property (3).

\medskip\noindent
Let $R$ be a Noetherian domain. Let $f : X \to Y$ be a morphism
of finite type schemes over $R$. By the discussion in the previous paragraph
it suffices to show that after replacing $R$ by $R_g$ for some $g \in R$
nonzero and $X$, $Y$ by their base changes to $R_g$,
formation of the scheme theoretic image of $f/R$ commutes
with arbitrary base change.

\medskip\noindent
Let $Y = V_1 \cup \ldots V_n$ be an affine open covering.
Let $U_i = f^{-1}(V_i)$. If the statement is true for each of the
morphisms $U_i \to V_i$ over $R$, then it holds for $f$.
Namely, the scheme theoretic image of $U_i \to V_i$
is the intersection of $V_i$ with the scheme theoretic
image of $f : X \to Y$ by Morphisms, Lemma
\ref{morphisms-lemma-quasi-compact-scheme-theoretic-image}.
Thus we may assume $Y$ is affine.

\medskip\noindent
Let $X = U_1 \cup \ldots U_n$ be an affine open covering.
Then the scheme theoretic image of $X \to Y$ is the same as
the scheme theoretic imge of $\coprod U_i \to Y$.
Thus we may assume $X$ is affine.

\medskip\noindent
Say $X = \Spec(A)$ and $Y = \Spec(B)$ and $f$ corresponds
to the $R$-algebra map $\varphi : A \to B$.
Then the scheme theoretic image of $f$ is $\Spec(A/\Ker(\varphi))$
and similarly after base change (by an affine morphism, but it is
enough to check for those).
Thus formation of the scheme theoretic image commutes
with base change if $\Ker(\varphi \otimes_R R') = \Ker(\varphi) \otimes_R R'$
for all ring maps $R \to R'$.

\medskip\noindent
After replacing $R$, $A$, $B$ by $R_g$, $A_g$, $B_g$ for a suitable
nonzero $g$ in $R$, we may assume $A$ and $B$ are flat over $R$.
By Lemma \ref{lemma-cokernel-flat}
we may also assume $B/A$ is a flat $R$-module.
Then $0 \to \Ker(\varphi) \to A \to B \to B/A \to 0$
is an exact sequence of flat $R$-modules, which implies
the desired base change statement.
\end{proof}